\section{USAMO 2005}
        \begin{enumerate}
        	\item (\textit{USAMO 2005}) (\textit{Zuming Feng}) Determine all composite positive integers $n$ for which it is possible to arrange all divisors of $n$ that are greater than 1 in a circle so that no two adjacent divisors are relatively prime.


 



\textbf{Solution 1 (official solution)}

No such circular arrangement exists for $n=pq$, where $p$ and $q$ are distinct primes. In that case, the numbers to be arranged are $p$; $q$ and $pq$, and in any circular arrangement, $p$ and $q$ will be adjacent. We claim that the desired circular arrangement exists in all other cases. If $n=p^e$ where $e\ge2$, an arbitrary circular arrangement works. Henceforth we assume that $n$ has prime factorization $p^{e_1}_{1}p^{e_2}_{2}\cdots p^{e_k}_k$, where $p_1<p_2<\cdots<p_k$ and either $k>2$ or else $\max(e1,e2)>1$. To construct the desired circular arrangement of $D_n:=\lbrace d:d|n\ \text{and}\ d>1\rbrace$, start with the circular arrangement of $n,p_{1}p_{2},p_{2}p_{3}\ldots,p_{k-1}p_{k}$ as shown.

 
                    \begin{figure}[H]
                        \centering
                        \includegraphics[width=0.5\textwidth]{Solution/USAMO_I_2005_1_Solution_Figure_1.png}
                    \end{figure}
                    

 Then between $n$ and $p_{1}p_{2}$, place (in arbitrary order) all other members of $D_n$ that have $p_1$ as their smallest prime factor. Between $p_{1}p_{2}$ and $p_{2}p_{3}$, place all members of $D_n$ other than $p_{2}p_{3}$ that have $p_2$ as their smallest prime factor. Continue in this way, ending by placing $p_k,p^{2}_{k},\ldots,p^{e_k}_{k}$ between $p_{k-1}p_k$ and $n$. It is easy to see that each element of $D_n$ is placed exactly one time, and any two adjacent elements have a common prime factor. Hence this arrangement has the desired property.

 \textit{Note.} In graph theory terms, this construction yields a Hamiltonian cycle in the graph with vertex set $D_n$ in which two vertices form an edge if the two corresponding numbers have a common prime factor. The graphs below illustrate the construction for the special cases $n=p^{2}q$ and $n=pqr$.

 


 



\textbf{Solution 2}

The proof that no arrangement exists for $n=pq$, where $p,q$ are distinct primes follows from above. Apply induction to prove all other cases are possible

 Base case: 

 
            \begin{enumerate}
                \item $n=p^r$, where $p$ is a prime and $r$ is a positive integer. Any arrangement suffices\item $n=pqr$, where $p,q,r$ are distinct primes. The following configuration works
            \end{enumerate}
            
 \[p,pq,pr,r,qr,q,pqr\]Inductive step: Suppose the desired arrangement exists for a composite $n$, show the arrangement exists for $np^r$, where $p$ is a prime relatively prime to $n$ and $r$ is a positive integer

 Let $a_1,a_2,\cdots,a_m$ be the arrangement of divisors of $n$, then $(a_i,a_{i+1})>1$ for $i=1,2,\cdots,m$, where $a_{m+1}=a_1$. The divisors of $np^r$ greater than 1 are of the form
 \[a_ip^j,p^j\qquad 1\leq i\leq m,1\leq j\leq r\]The following sequence works
 \[a_1,\cdots,a_{m-1}, a_{m-1}p,\text{other divisors in arbitrary order},a_mp,a_m\]since all other divisors are divisible by $p$.

 <i>Alternate solutions are always welcome. If you have a different, elegant solution to this problem, please \href{https://artofproblemsolving.com/wiki/index.php?title=2005\_USAMO\_Problems/Problem\_1&action=edit}{add it} to this page.</i>

 


	\item (\textit{USAMO 2005}) (\textit{Răzvan Gelca}) Prove that the
system

 \begin{align*} x^6+x^3+x^3y+y & = 147^{157} \\ x^3+x^3y+y^2+y+z^9 & = 157^{147} \end{align*}

has no solutions in integers $x$, $y$, and $z$.


 



\textbf{Solution 1}

It suffices to show that there are no solutions to this system in the integers mod 19.  We note that $152 = 8 \cdot 19$, so $157 \equiv -147 \equiv 5 \pmod{19}$.  For reference, we construct a table of powers of five:
 \[\begin{array}{c|c||c|c} n& 5^n &n & 5^n \\\hline 1 & 5 & 6 & 7 \\ 2 & 6 & 7 & -3 \\ 3 & -8 & 8 & 4 \\ 4 & -2 & 9 & 1 \\ 5 & 9 && \end{array}\]Evidently, the order of 5 is 9.  Hence 5 is the square of a multiplicative generator of the nonzero integers mod 19, so this table shows all nonzero squares mod 19, as well.

 It follows that $147^{157} \equiv (-5)^{13} \equiv -5^4 \equiv 2$, and $157^{147} \equiv 5^3 \equiv -8$.  Thus we rewrite our system thus:
 \begin{align*} (x^3+y)(x^3+1) &\equiv 2 \\ (x^3+y)(y+1) + z^9 &\equiv -8 . \end{align*}
Adding these, we have
 \[(x^3+y+1)^2 - 1 + z^9 \equiv -6,\]or
 \[(x^3+y+1)^2 \equiv -z^9 - 5 .\]By \href{https://artofproblemsolving.com/wiki/index.php?title=Fermat\%27s\_Little\_Theorem}{Fermat's Little Theorem}, the only possible values of $z^9$ are $\pm 1$ and 0, so the only possible values of $(x^3+y+1)^2$ are $-4,-5$, and $-6$.  But none of these are squares mod 19, a contradiction.  Therefore the system has no solutions in the integers mod 19.  Therefore the solution has no equation in the integers.  $\blacksquare$

 



\textbf{Solution 2}

Note that the given can be rewritten as 

 (1) $(x^3+y)(x^3+1) = (x^3+y)(x+1)(x^2-x+1) = 147^{157} = 7^{314}3^{152}$,

 (2) $(x^3+y)(y+1)+z^9 = 157^{147} \rightarrow (x^3+y)(y+1) = (157^{49}-z^3)(157^{98}+157^{49}z^3+z^6)$.

 We can also see that

 $x^2-x+1 = (x+1)^2-3(x+1)+3 \rightarrow \gcd(x+1, x^2-x+1) \le 3$.

 Now we notice 

 $x^3+1 = 3^a7^b$

 for some pair of non-negative integers $(a,b)$. We also note that 

 $x^2 \ge 2x$ when $|x| \ge 2 \rightarrow x^2-x+1 \ge x+1$ 

 when $|x| \ge 2$. If $x = 1$ or $x = -1$ then examining (1) would yield $147^{157} \equiv 0 \pmod{2}$ which is a contradiction. If $x = 0$ then from (1) we can see that $y+1 = 147^{157}$, plugging this into 2 yields 

 (3) $(147^{157}-1)(147^{157}) = (157^{49}-z^3)(157^{98}+157^{49}z^3+z^6)$, $146 | (147^{157}-1)$, $146 = 2*73$.

 Noting that 73 is a prime number we see that it must divide at least 1 of the 2 factors on the right hand side of 3. Let us consider both cases.

 $73 | (157^{49}-z^3) \rightarrow z^3 \equiv 11^{49} \pmod{73} \rightarrow (11^{49})^{24} \equiv 1 \pmod{73}$.

 However

 $(11^{49})^{24} \equiv 8 \pmod{73}$

 Thus we can see that 73 cannot divide the first factor in the right hand side of (3). Let us consider the next case.

 $73 | (157^{98}+157^{49}z^3+z^6) \rightarrow 11^{98}+11^{49}z^3+z^6 \equiv 0 \pmod{73}$.

 However

 $11^{98}+11^{49}z^3+z^6 \equiv z^6+47z^3+19 \equiv z^6-26z^3+19 \equiv (z^3-13)^2-150 \equiv (z^3-13)^2-4 \equiv (z^3-11)(z^3-15) \pmod{73}$.

 It can be seen that 11 and 15 are not perfect cubes $\pmod{73}$ from the following.

 $15^{24} \equiv 11^{24} \equiv 8 \not\equiv 1 \pmod{73}$

 We can now see that $|x| \ge 2$. Furthermore, notice that if

 $x+1 = \pm3^k7^j$

 for a pair of positive integers $(j,k)$ means that 

 $|x^2-x+1| \le 3$

 which cannot be true. We now know that 

 $x+1 = \pm3^k, \pm7^k \rightarrow x^2-x+1 = 3*7^m, 3^m$.

 Suppose that 

 $x+1 = \pm7^k \rightarrow (x+1)^2-3(x+1)+3 = 3^m \rightarrow 7^{2k}\pm3*7^k+3  = 3^m \rightarrow 7^{2k} \equiv 0\pmod{3}$ 

 which is a contradiction. Now suppose that 

 $x+1 = \pm3^k \rightarrow (x+1)^2-3(x+1)+3 = 3*7^m \rightarrow 3^{2k}\pm3^{k+1}+3 = 3*7^m \rightarrow 3^{2k-1}\pm3^{k}+1 = 7^m \rightarrow 3^k(3^{k-1}\pm1) = 7^m-1$.

 We now apply the lifting the exponent lemma to examine the power of 3 that divides each side of the equation when $m > 0$ to obtain

 $k \le 2+v_3(m) \le 2+log_3(m) \rightarrow 7^m-1 = 3^k(3^{k-1}\pm1)\le 3^{2+log_3(m)}(3^{1+log_3(m)}\pm1) = 9m(3m\pm1)$.

 For $m > 2$ we can see that 

 $7^m-1 > 9m(3m\pm1)$ 

 which is a contradiction. Therefore there only possible solution is when 

 $m = 0 \rightarrow k = 1 \rightarrow x = 2, m = 1 \rightarrow k = 1 \rightarrow x = -4, m = 2 \rightarrow$ no integer solutions for k.

 Plugging this back into (1) and (2) yields

 (4) $3^{155} | (x^3+y) \rightarrow 3^{155}| (157^{49}-z^3)(157^{98}+157^{49}z^3+z^6)$.

 In order for (4) to be true we must have 9 dividing at least 1 of the factors on the right hand side of the equation. Let us consider both cases.

 $9 | 157^{49}-z^3 \rightarrow 157^{49}-z^3 \equiv 0 \pmod{9}$.

 However,

 $z^3 \equiv 0, 1, 8\pmod{9}, 157^{49}-z^3 \equiv 4-z^3 \not\equiv 0 \pmod{9}$.

 We now consider the second case.

 $9 | (z^6+157^{49}z^3+157^{98}) \rightarrow 157^{98}+157^{49}z^3+z^6\equiv 0\pmod{9}$.

 However

 $z^6+157^{49}z^3+157^{98} \equiv z^6+4z^3+7 \equiv (z^3+2)^2+3 \not\equiv 0\pmod{9}$

 Therefore there are no solutions to the given system of diophantine equations. $\blacksquare$

 



\textbf{Solution 3}

We will show there is no solution to the system modulo 13. Add the two equations and add 1 to obtain
 \[(x^3 + y + 1)^2 + z^9 = 147^{157} + 157^{147} + 1.\]By Fermat's Theorem, $a^{12}\equiv 1\pmod{13}$ when $a$ is not a multiple of 13. Hence we compute $147^{157}\equiv 4^1\equiv 4\pmod{13}$ and $157^{147}\equiv 1^3\equiv 1\pmod{13}$. Thus
 \[(x^3 + y + 1)^2 + z^9\equiv 6\pmod{13}.\]The cubes mod 13 are $0$, $\pm 1$, and $\pm 5$. Writing the first equation as
 \[(x^3 + 1)(x^3 + y)\equiv 4\pmod{13},\]we see that there is no solution in case $x^3\equiv -1\pmod{13}$ and for $x^3$ congruent to $0, 1, 5, -5\pmod{13}$, correspondingly $x^3 + y$ must be congruent to $4, 2, 5, -1$. Hence
 \[(x^3 + y + 1)^2\equiv \text{12, 9, 10, or 0}\pmod{13}.\]Also $z^9$ is a cube, hence $z^9$ must be $\text{0, 1, 5, 8, or 12}\pmod{13}$. It is easy to check that $6\pmod{13}$ is not obtained by adding one of $0, 9, 10, 12$ to one of $0, 1, 5, 8, 12$. Hence the system has no solutions in integers.

 \textit{Note}: This argument shows there is no solution even if $z^9$ is replaced by $z^3$.

 <i>Alternate solutions are always welcome. If you have a different, elegant solution to this problem, please \href{https://artofproblemsolving.com/wiki/index.php?title=2005\_USAMO\_Problems/Problem\_2&action=edit}{add it} to this page.</i>

 


	\item (\textit{USAMO 2005}) (\textit{Zuming Feng}) Let $ABC$ be an acute-angled triangle, and let $P$ and $Q$ be two points on side $BC$. Construct point $C_1$ in such a way that convex quadrilateral $APBC_1$ is cyclic, $QC_1 \parallel CA$, and $C_1$ and $Q$ lie on opposite sides of line $AB$. Construct point $B_1$ in such a way that convex quadrilateral $APCB_1$ is cyclic, $QB_1 \parallel BA$, and $B_1$ and $Q$  lie on opposite sides of line $AC$.  Prove that points $B_1, C_1,P$, and $Q$ lie on a circle.


 



\textbf{Solution 1}

Let $B_1'$ be the second intersection of the line $C_1A$ with the circumcircle of $APC$, and let $Q'$ be the second intersection of the circumcircle of $B_1' C_1P$ and line $BC$.  It is enough to show that $B_1'=B_1$ and $Q' =Q$.  All our angles will be directed, and measured mod $\pi$.

 
                    \begin{figure}[H]
                        \centering
                        \includegraphics[width=0.5\textwidth]{Solution/USAMO_I_2005_3_Solution_Figure_1.png}
                    \end{figure}
                    

 Since points $C_1, P, Q', B_1'$ are concyclic and points $C_1, A,B_1'$ are collinear, it follows that
 \[\angle C_1 Q' P \equiv \angle C_1 B_1' P \equiv \angle A B_1' P .\]But since points $A, B_1', P, C$ are concyclic,
 \[\angle AB_1'P \equiv \angle ACP .\]It follows that lines $AC$ and $C_1 Q'$ are parallel.  If we exchange $C$ with $B$ and $C_1$ with $B_1'$ in this argument, we see that lines $AB$ and $B_1' Q'$ are likewise parallel.

 It follows that $Q'$ is the intersection of $BC$ and the line parallel to $AC$ and passing through $C_1$.  Hence $Q' = Q$.  Then $B_1'$ is the second intersection of the circumcircle of $APC$ and the line parallel to $AB$ passing through $Q$.  Hence $B_1' = B_1$, as desired.  $\blacksquare$

 Motivation: One can notice that if you take such a $B_1'$ then $PQB_1'C_1$ is cyclic, and that similarly $PQB_1C_1'$ is also cyclic. One gets the intuition that only one such circle should exist where the other chord passes through A, and so sets up a ghost point for $Q$, which works. ~cocohearts

 



\textbf{Solution 2}

\textit{Lemma}. $B_1, A, C_1$ are collinear.

 Suppose they are not collinear. Let line $B_1 A$ intersect circle $ABP$ (i.e. the circumcircle of $ABP$) again at $C_2$ distinct from $C_1$. Because $\angle C_2 B_1 Q = \angle C_2 AB = \angle C_2 PB = 180^\circ - \angle C_2 PQ$, we have that $B_1 C_2 PQ$ is cyclic. Hence $\angle C_2 QP = \angle C_2 B_1 P = \angle P B_1 A = \angle C$, so $C_2 Q // AC$. Then $C_2$ must be the other intersection of the parallel to $AC$ through $Q$ with circle $ABP$. Then $C_2$ is on segment $C_1 Q$, so $C_2$ is contained in triangle $ABQ$. However, line $AB_1$ intersects this triangle only at point $A$ because $B_1$ lies on arc $AC$ not containing $P$ of circle $APC$, a contradiction. Hence, $B_1, A, C_1$ are collinear, as desired.

 As a result, we have $\angle C_1 B_1 Q = \angle C_1 AB = \angle C_1 PB = 180^\circ - \angle C_1 PQ$, so $B_1 C_1 PQ$ is cyclic, as desired.

 



\textbf{Solution 3}

Due to the parallel lines, $m\angle C_1QB_1=m\angle A.$ Therefore, it suffices to prove that 
 \[m\angle C_1PB_1=m\angle C_1BA+m\angle ACB_1=m\angle A.\] Note that $m\angle BC_1A+m\angle AB_1C=180^{\circ}$ by the cyclic quadrilaterals. Now, the condition simplifies to proving $C_1,A,B_1$ collinear.

 Use barycentric coordinates. Let 
 \[A=(1,0,0), B=(0,1,0), C=(0,0,1), P=(0,p,1-p), Q=(0,q,1-q).\] By the parallel lines, 
 \[C_1=(1-q-z',q,z'), B_1=(q-y',y',1-q)\] for some $y',z'.$ As $A$ is a vertex of the reference triangle, we must prove that $q(1-q)=y'z'.$

 Now we find the circumcircle of $\triangle APC.$ Let its equation be 
 \[-a^2yz-b^2xz-c^2xy+(x+y+z)(ux+vy+wz)=0.\] Substituting in $A,C$ gives $u=w=0.$ Substituting in $P$ yields 
 \[-a^2p(1-p)+vp=0\implies v = a^2(1-p).\] Substituting in $B_1$ yields
 \[0=-a^2y'(1-q)-b^2(q-y')(1-q)-c^2y'(q-y')+a^2(1-p)y'=0.\]
 \[c^2y'^2+(a^2q-a^2p+b^2(1-q)-c^2q)y'-b^2q(1-q)=0\]

 Similarly, the circumcircle of $\triangle APB$ is 
 \[-a^2yz-b^2xz-c^2xy+(x+y+z)a^2pz.\] Now, we substitute $C_1$ into the equation

 
 \[0=-a^2qz'-b^2z'(1-q-z')-c^2q(1-q-z')+a^2pz'=0\]
 \[b^2 z'^2+(a^2p-a^2q-b^2(1-q)+c^2q)z'-c^2q(1-q)=0\]

 Let $k=a^2p-a^2q-b^2(1-q)+c^2q,$ and note that both $y',z'$ are negative. By the quadratic formula, we have
 \[y'=\frac{k-\sqrt{k^2+4b^2c^2q(1-q)}}{2c^2}, z'=\frac{-k-\sqrt{k^2+4b^2c^2q(1-q)}}{2b^2}\]Multiplying these two, we have 
 \[y'z'=\frac{4b^2c^2q(1-q)}{4b^2c^2}=q(1-q),\] as desired. $\blacksquare$

 <i>Alternate solutions are always welcome. If you have a different, elegant solution to this problem, please \href{https://artofproblemsolving.com/wiki/index.php?title=2005\_USAMO\_Problems/Problem\_3&action=edit}{add it} to this page.</i>

 


	\item (\textit{USAMO 2005}) Legs $L_1,L_2,L_3,L_4$ of a square table each have length $n$, where $n$ is a positive integer. For how many ordered 4-tuples $\left(k_1,k_2,k_3,k_4\right)$ of nonnegative integers can we cut a piece of length $k_i$ from the end of leg $L_i\ (i=1,2,3,4)$ and still have a stable table?


 (The table is stable if it can be placed so that all four of the leg ends touch the floor. Note that a cut leg of length 0 is permitted.)


 



\textbf{Solution 1}

First flip the table upside down and observe the plane that the ends of the legs form. Let this plane be $ax+by+c=z$. Assume WLOG that the table is the square with endpoints $(0,0), (0,1), (1,0), (1,1)$ on the xy plane. Thus the lengths of the legs will be $c, c+b, c+a, c+a+b$ respectively. In other words $k_1+k_3=k_2+k_4=x$ for some $0 \le x \le 2n$.

 For $0 \le x \le n$, it is easy to see that the number of stable tables is $(x+1)^2$.

 For $n+1 \le x \le 2n$, we can biject the stable tables of $x$ to the stable of $2n-x$ by replacing all of the legs with the cut off portions.

 So the total number of stable tables is $1^2+2^2+...+n^2+(n+1)^2+n^2+...+2^2+1^2 =$ $(1^2+2^2+...+(n+1)^2)+(1^2+2^2+...+n^2) =$ $\dfrac{(n+1)(n+2)(2n+3)}{6}+\dfrac{n(n+1)(2n+1)}{6}=\boxed{\dfrac{2n^3+6n^2+7n+3}{3}}$.

 


 



\textbf{Solution 2}

Let $C_i$ be the number of sets of cuts $\{k_1,k_2,k_3,k_4\}$ that result in the longest leg being of length $i$.  Then $C_i=C_{i-1}+x$, and we will evaluate $x$ to develop a recursion.  Note that if a set of cuts has non-zero cuts, then if that set is counted in $C_i$, subtract $1$ from each of the cuts to obtain a set of cuts that is counted in $C_{i-1}$.  For example, if $\{2,3,4,5\}$ was counted in $C_5$, then $\{1,2,3,4\}$ would have been counted in $C_4$ because subtracting the same length from each leg preserves the quality of the legs being coplanar, and there is a bijection.

 Now we must evaluate those sets of cuts that don't fall in this bijection, which are clearly those where one leg is completely cut off.  After some simple geometric reasoning that I won't include here, we see that one corner has a leg of length of zero, and the opposite leg, called the initial leg, will have length $i$  If the remaining legs have length $y$ and $z$, then $(y,z)=(0,i),(1,i-1),...,(i-1,1),(i,0)$, so $i+1$ options.  The initial corner can be any of the four legs, but each of the four permutations of the set of cuts ${0,0,i,i}$ is repeated twice, so we have $4(i+1)-4=4i$ total additional sets of cuts.  We conclude that $C_i=C_{i-1}+4i$, and note that $C_0=1$.  Now we can create generating functions.  $F(x)=\sum_{i=0}^\infty C_ix^i$.  Also, $G(x)=\sum_{i=1}^\infty 4ix^i$. From the recursion, we have
 \[F(x)=C_0+xF(x)+G(x)\Longrightarrow F(x)=\frac{1+G(x)}{1-x}\]This final equation is easy to analyze.  Simply use $G(x)$ as the first term of a geometric series with constant multiplier of $x$.  This gives $C_i=2i^2+2i+1$.  The total number of sets of cuts for a table of legs of length $n$ is the sum of all the $C_i$, $0\le i\le n$, from which we deduce the answer $\frac{2n^3 + 6n^2 +7n + 3}{3}$, or $\frac{(n+1)(2n^2 + 4n +3)}{3}$.

 


 



\textbf{Solution 3}

The problem's condition is that the ends of the legs must be coplanar. So the plane formed by legs $L_1$, $L_2$, and $L_3$ is the same as the plane formed by the legs $L_2$, $L_3$, and $L_4$ iff the problem's condition is met.

 We now assign coordinates to the ends of the legs. Let the table be on the xy-plane such that $L_1$ is at $(0,0,k_1)$, $L_2$ is at $(s,0,k_2)$, $L_3$ is at $(0,s,k_3)$, and $L_4$ is at $(s,s,k_4)$, where $s$ is the side length of the table. (Note that we have turned the table over so that the top lies on the xy plane while the legs point in the +z direction - this does not affect the result.) A plane in 3-space has the equation form $Ax+By+Cz=1$. Plugging the coordinates of $L_1$, $L_2$, and $L_3$ into $Ax+By+Cz=1$, we eventually get $A_1=\dfrac{k_1-k_2}{sk_1}$, $B_1=\dfrac{k_1-k_3}{sk_1}$, $C_1=\dfrac{1}{k_1}$. Plugging the coordinates of $L_2$, $L_3$, and $L_4$ into $Ax+By+Cz=1$, we get $A_2=\dfrac{k_3-k_4}{s(k_2+k_3-k_4)}$, $B_2=\dfrac{k_2-k_4}{s(k_2+k_3-k_4)}$, $C_2=\dfrac{1}{k_2+k_3-k_4}$. 

 Since we need these to be coplanar, we set $A_1=A_2$, $B_1=B_2$, and $C_1=C_2$. (Notice that the $s$'s cancel.) From $A_1=A_2$, we get $(k_1-k_2)(k_2+k_3-k_4)=k_1(k_3-k_4)$ [1]. From $B_1=B_2$, we get $(k_1-k_3)(k_2+k_3-k_4)=k_1(k_2-k_4)$ [2]. From $C_1=C_2$, we get $k_1=k_2+k_3-k_4$ [3]. Notice that [1] and [2] can both be obtained directly from [3] and are therefore redundant. The only condition we have left to ensure that the ends of the legs are coplanar is $k_1+k_4=k_2+k_3$. Let $k_1+k_4=k_2+k_3=m$.

 We proceed using casework by $m$. $m$ is an integer between $0$ and $2n$, inclusive. Furthermore, we note that $k_i$ range from $0$ to $n$, inclusive. For $m=0$, each of $(k_1,k_4)$ and $(k_2,k_3)$ have one choice, $(0,0)$. For $m=1$, they each have two choices, $(0,1)$ and $(1,0)$. And so on, until $m=n$, where they each have $n+1$ choices: $(0,n),(1,n-1),\cdots,(n-1,1),(n,0)$. When $m=n+1$, the restriction that $0\le k_i\le n$ takes effect: there are only $n$ choices for each ordered pair, which are $(1,n),(2,n-1),\cdots,(n-1,2),(n,1)$. And so on, until $m=2n$, where there is one choice for each ordered pair, $(n,n)$.

 The number of possible 4-tuples $(k_1,k_2,k_3,k_4)$ is thus $1^2+2^2+...+n^2+(n+1)^2+n^2+...+2^2+1^2$, which is equivalent to $(1^2+2^2+...+(n+1)^2)+(1^2+2^2+...+n^2)$. Our final answer is $\dfrac{(n+1)(n+2)(2n+3)}{6}+\dfrac{n(n+1)(2n+1)}{6}=\boxed{\dfrac{2n^3+6n^2+7n+3}{3}}$.

 
<i>Alternate solutions are always welcome. If you have a different, elegant solution to this problem, please \href{https://artofproblemsolving.com/wiki/index.php?title=2005\_USAMO\_Problems/Problem\_4&action=edit}{add it} to this page.</i>

 


	\item (\textit{USAMO 2005}) Let $n$ be an integer greater than 1. Suppose $2n$ points are given in the plane, no three of which are collinear. Suppose $n$ of the given $2n$ points are colored blue and the other $n$ colored red. A line in the plane is called a \textit{balancing line} if it passes through one blue and one red point and, for each side of the line, the number of blue points on that side is equal to the number of red points on the same side.


 Prove that there exist at least two balancing lines.


 



\textbf{Solution 1}

Consider the convex hull of the the $2n$ points, or the points that would form the largest convex polygon. If the points in the convex hull contain both red and blue points, then there must be at least 2 edges of the graph of the convex hull such that the edge connects a blue and a red point. Drawing a line through those points would give a balancing line, as we have n-1 blue points and n-1 red points on one side, and 0 points on the other.

 Therefore it suffices to show that there exist at least 2 balancing lines when the convex hull is colored all the same color. 

 Pick a random point on the convex hull, and without loss of generality we can say it is blue (if there are no red we can change all the colors, and end up with an equivalent setup). Consider a line going through it and not any other points. As we rotate the line clockwise, we encounter the red points in some order. Let the ith point encountered be $R_i$. Let $b_i$ and $r_i$ be the number of points encountered before $R_i$. Then $r_i=i-1$.

 Define a sequence $f(i)=b_i-r_i$. Then $f(1)=b_i-(1-1)=b_i\geq0$$f(n)=b_n-(n-1)=b_n-n+1\leq0$, because we can only encounter up to n-1 blue points.Thus, $f(i)$ goes from negative to positive as $i$ goes from 1 to $n$. We can also see that $f(i)$ can only increase by one for each change in i, so we know $f(i)$ must be 0 for some value of $i$. Take the first $f(i)$ where $f(i) = 0$. Since $f(i-1)$ must have been positive, the ith point must have been red. Thus there is a balancing line for every point on the convex hull. Since a polygon must have at least 3 vertices, there must be at least $3>2$ balancing lines for the set of points when the convex hull is all the same color, and the statement is true as we desired.

 <i>Alternate solutions are always welcome. If you have a different, elegant solution to this problem, please \href{https://artofproblemsolving.com/wiki/index.php?title=2005\_USAMO\_Problems/Problem\_5&action=edit}{add it} to this page.</i>

 



\textbf{Solution 2}

Case 1 One can notice That if any blue and red point have not points on one of it's sides this mean the line between them is a balancing line (as on the other side we will have $n-1$ blue points and $n-1$ red).

 Case 2: Now we consider the case in which all the points are inside a polygon with all vertices being one color (assume WLOG it's blue

 Now consider a line that is pivoting at any vertex and rotating till reaching another vertex of the polygon (convex hull) passing through every point in side the polygon. and let $x$ be the number of red points.minus the number of blue points on one side of the line.In other words every time the line passes through a blue point $x$ will increase by $1$, and decrease by $1$ if it passes through a red point.Firstly we have that $x=1$ and when the line reaches the other vertex of the polygon  it will be equal to $-1$.So when $x$ will be equal to zero, we have a balancing line.Now repeating this process on all the other vertices of the polygon (convex hull) then we will have balancing lines equal to the number of the vertices.

 By Ywgh1

 


	\item (\textit{USAMO 2005}) For $m$ a positive integer, let $s(m)$ be the sum of the digits of $m$. For $n\ge 2$, let $f(n)$ be the minimal $k$ for which there  exists a set $S$ of $n$ positive integers such that  $s\left(\sum_{x\in X} x\right) = k$ for any nonempty subset $X\subset S$.  Prove that there are constants $0 < C_1 < C_2$ with

 \[C_1 \log_{10} n \le f(n) \le C_2 \log_{10} n.\]


 



\textbf{Solution}

For the upper bound, let $p$ be the smallest integer such that $10^p\geq n(n+1)/2$ and let
 \[S = \{10^p - 1, 2(10^p - 1), \ldots, n(10^p - 1)\}.\]The sum of any nonempty set of elements of $S$ will have the form $k(10^p - 1)$ for some $1\leq k\leq n(n+1)/2$. Write $k(10^p - 1) = [(k - 1)10^p] + [(10^p - 1) - (k - 1)]$. The second term gives the bottom $p$ digits of the sum and the first term gives at most $p$ top digits. Since the sum of a digit of the second term and the corresponding digit of $k - 1$ is always 9, the sum of the digits will be $9p$. Since $10^{p-1} < n(n+1)/2$, this example shows that
 \[f(n)\leq 9p\leq 9\log_{10}(5n(n+1)).\]Since $n\geq 2$, $5(n + 1) < n^4$, and hence
 \[f(n) < 9\log_{10}n^5 = 45\log_{10}n.\]For the lower bound, let $S$ be a set of $n\geq 2$ positive integers such that any nonempty $X\subset S$ has $s\left(\sum_{x\in X} x\right) = f(n)$. Since $s(m)$ is always congruent to $m$ modulo 9, $\sum_{x\in X} x\equiv f(n)\pmod{9}$ for all nonempty $X\subset S$. Hence every element of $S$ must be a multiple of 9 and $f(n)\geq 9$. Let $q$ be the largest positive integer such that $10^q - 1\leq n$. Lemma 1 below shows that there is a nonempty subset $X$ of $S$ with $\sum_{x\in X} x$ a multiple of $10^q - 1$, and hence Lemma 2 shows that $f(n)\geq 9q$.

 \textbf{Lemma 1.} Any set of $m$ positive integers contains a nonempty subset whose sum is a multiple of $m$.

 \textit{Proof.} Suppose a set $T$ has no nonempty subset with sum divisible by $m$. Look at the possible sums mod $m$ of nonempty subsets of $T$. Adding a new element $a$ to $T$ will give at least one new sum mod $m$, namely the least multiple of $a$ which does not already occur. Therefore the set $T$ has at least $|T|$ distinct sums mod $m$ of nonempty subsets and $|T| < m$. $\blacksquare$

 \textbf{Lemma 2.} Any positive multiple $M$ of $10^q - 1$ has $s(M)\geq 9q$.

 \textit{Proof.} Suppose on the contrary that $M$ is the smallest positive multiple of $10^q - 1$ with $s(M) < 9q$. Then $M\neq 10^q - 1$, hence $M > 10^q$. Suppose the most significant digit of $M$ is the $10^m$ digit, $m\geq q$. Then $N = M - 10^{m-q}(10^q - 1)$ is a smaller positive multiple of $10^q - 1$ and has $s(N)\leq s(M) < 9q$, a contradiction. $\blacksquare$

 Finally, since $10^{q+1} > n$, we have $q + 1 > \log_{10} n$. Since $f(n)\geq 9q$ and $f(n)\geq 9$, we have
 \[f(n)\geq \frac{9q + 9}{2} > \frac{9}{2}\log_{10}n.\]Weaker versions of Lemmas 1 and 2 are still sufficient to prove the desired type of lower bound.

 <i>Alternate solutions are always welcome. If you have a different, elegant solution to this problem, please \href{https://artofproblemsolving.com/wiki/index.php?title=2005\_USAMO\_Problems/Problem\_6&action=edit}{add it} to this page.</i>-mahthir marshad

 


\end{enumerate}